\section{The first two derivative regimes}
\label{sec:jacobian}

We prove both main theorems when $p\nmid d-1$. The common
normalization and coefficient lemmas hold in every characteristic.
We first treat the nonzero Jacobian scalar, then the
derivative-zero branch separately.

\subsection{Adjacent Jacobians when \texorpdfstring{$p\nmid d(d-1)$}{p does not divide d(d-1)}}

Fix $M\ge1$, $\deg h\le M$, and $n=n_J$. Suppose
\begin{equation}
 F_j\mid F_j'H_j(h),\qquad j=n,n+1.
 \label{eq:two-jacobians}
\end{equation}
Write $h=\Deltaf Q+v$ by \cref{lem:normal}. Telescoping in
\eqref{eq:telescoping-jacobian} shows that
\eqref{eq:two-jacobians} is unchanged upon replacing $h$ by $v$,
and $\deg v\le M$.

Suppose $v$ has positive leading degree $D$ and leading
coefficient $a_D\ne0$. For the target $M_D$ of
\cref{lem:leading}, put
\[
 C_j=[M_D](\Pi_j\cH_j(v)).
\]
Both selected levels exceed $2\lfloor\log_dD\rfloor$ and meet
the stabilization bound for the cap $M$. Equations
\eqref{eq:leading-coefficient}, \eqref{eq:telescoping-jacobian},
and \eqref{eq:two-jacobians} yield
\begin{equation}
 d^nC_n=a_D,\qquad d^{n+1}C_{n+1}=a_D.
 \label{eq:two-leading-jacobians}
\end{equation}
By \cref{prop:stabilization}, $C_{n+1}=C_n$. Subtract $d$
times the first equality from the second to get
$(1-d)a_D=0$. The hypothesis $p\nmid d-1$ contradicts
$a_D\ne0$.

Thus $v=a_0$ is constant. At return $j$ the remaining condition
is $F_j\mid ja_0F_j'$. Here $F_j'$ has nonzero leading
coefficient $d^j$ and degree $d^j-1<\deg F_j$, so
$ja_0=0$. Applying this at $j=n,n+1$ and subtracting gives
$a_0=0$. The two Jacobian tests therefore imply $h\in\Bf$.

\subsection{Squarefree returns when \texorpdfstring{$p\mid d$}{p divides d}}

In this case $f'=0$, and hence
\begin{equation}
 F_j'=-1\qquad(j\ge1).
 \label{eq:derivative-zero}
\end{equation}
Every $F_j$ is squarefree. The tests in the second row of
\cref{tab:tests} give $\cH_n(v)=\cH_{n+1}(v)=0$ after
normalization. A positive normal degree $D\le M$ is excluded
at either level by \cref{lem:leading}. For a constant normal
part, the two equations are $na_0=0$ and $(n+1)a_0=0$,
and again $a_0=0$. This argument includes $p=2$, for which
$-1=1$; it does not replace the iteration of $f$ by another clock.

\subsection{Finite equivalence and ordinary detection}

If $h=\Deltaf Q$, then
\begin{equation}
 H_j(h)=Q(f^{\circ j}(x))-Q(x)
\quad\text{is divisible by }F_j
 \label{eq:polynomial-telescoping}
\end{equation}
for every $j$. It satisfies both the applicable divisibilities
and the ordinary-root tests. In the first regime, ordinary-root
vanishing implies the Jacobian test by
\cref{lem:root-certificates}; in the second it is equivalent
to divisibility because of \eqref{eq:derivative-zero}.
Together with the two preceding subsections, this proves the
three-way equivalence in \cref{thm:finite} whenever $p\nmid d-1$.

If $h\in\Kf$, formula \eqref{eq:ordinary-repetition} supplies
the two ordinary-root tests, so $h\in\Bf$. The reverse inclusion
is telescoping. This proves \cref{thm:rigidity} in these regimes.
Alternatively, after excluding every positive normal part, an
ordinary fixed point of $f$ directly excludes a nonzero constant:
the polynomial $f(x)-x$ has a root in the algebraically closed
field $k$.

For a noncoboundary, one ordinary-root test fails at a point $a$
of period $r\mid j$, where $j=n$ or $n+1$. Then
\[
 H_j(h)(a)=(j/r)S_h(O_a)\ne0
\]
implies $S_h(O_a)\ne0$, without division in $k$. In particular
the detecting primitive period is at most $n+1$. With
$m=\lfloor\log_dM\rfloor$,
\[
 \deg F_j\le d^{n+1}=d^{3m+5}\le d^5M^3.
\]
This is the first bound in \eqref{eq:degree-bounds}, including
zero and constant inputs in the equivalence proved above.
